% Computer Science A --- \texttt{java.lang.Class}, reflection (\texttt{java.lang.reflect}),
% and concurrency (\texttt{java.lang.Thread} / \texttt{Runnable}).
% Stream API questions live in \texttt{cs-a/functional.tex} (after lambdas/method references); stream pipelines tie to parallel execution (\texttt{parallelStream}, fork/join) alongside threads here.
% Reflection and thread \emph{free-response} prompts migrated from \texttt{cs-a/03-practice.tex} sit just before the matching MCQ blocks below.
% Order: class metadata, reflection types, threads.
% Heap/stack basics: \texttt{cs-a/lang.tex}; streams + lambdas: \texttt{cs-a/functional.tex}.

\runningheader{Computer Science A}{Reflection, class loading, \& threads}{Page \thepage\ of \numpages}

\begingradingrange{csajavaplus}

\uplevel{\par\textbf{Class objects and dynamic loading} (\texttt{java.lang.Class})\par}

\question[4]
Retrieving Classes. Which is the correct way to retrieve references to classes in Java?
\begin{choices}
  \choice \texttt{obj.getClass()} method
  \choice \texttt{Object.class}
  \choice \texttt{Class.forName()} method
  \CorrectChoice All of the above
\end{choices}
\begin{solution}
  All three yield a \texttt{Class<?>} (or a concrete type parameter) in different situations:

\begin{lstlisting}[style=javaexam]
Object obj = "hello";
Class<?> c1 = obj.getClass();          // runtime class of instance

Class<Object> c2 = Object.class;        // class literal (compile-time)

Class<?> c3 = Class.forName("java.lang.String"); // load by name (may throw)
\end{lstlisting}

  Literals are resolved at compile time; \texttt{forName} is used for plugins, JDBC drivers, serialization, frameworks, etc.
\end{solution}

\question[4]
\lstinline!Class.forName("SuperLaser")! returns an \textbf{instance} of class \texttt{SuperLaser}.
\begin{choices}
  \CorrectChoice False
  \choice True
\end{choices}
\begin{solution}
  \texttt{Class.forName(String)} returns a \texttt{Class<?>} representing the type---metadata, not \texttt{new SuperLaser()}. Creating an instance is a separate step (constructors, DI, factories):

\begin{lstlisting}[style=javaexam]
Class<?> c = Class.forName("com.example.SuperLaser");
Object instance = c.getDeclaredConstructor().newInstance();
\end{lstlisting}

  \texttt{Class.newInstance()} existed in older Java but is deprecated; prefer \texttt{getDeclaredConstructor(...).newInstance(...)} with accessible checks as needed.
\end{solution}


\newpage

\uplevel{\par\textbf{Reflection} (\texttt{java.lang.reflect})\par}

\question[4]
Reflection API. Which of these is \textbf{NOT} a class in the Reflection API in Java?
\begin{choices}
  \choice \texttt{Constructor}
  \choice \texttt{Field}
  \choice \texttt{Method}
  \choice \texttt{Modifier}
  \CorrectChoice \texttt{Variable}
\end{choices}
\begin{solution}
  Core reflection types live in \textbf{\texttt{java.lang.reflect}} and on \texttt{java.lang.Class}. Typical imports:

\begin{lstlisting}[style=javaexam]
import java.lang.reflect.Constructor;
import java.lang.reflect.Field;
import java.lang.reflect.Method;
import java.lang.reflect.Modifier;
\end{lstlisting}

  \texttt{Class}, \texttt{Constructor}, \texttt{Method}, and \texttt{Field} describe types and members loaded from bytecode. \texttt{Modifier} is a utility class for decoding \texttt{int} modifier bitmasks. There is \textbf{no} \texttt{java.lang.reflect.Variable} type used the same way: local variables are not first-class reflection objects; only \texttt{Field} represents declared fields on a class.
\end{solution}

\question[4]
Which type is \emph{not} part of the core Java reflection types used to inspect classes and members?
\begin{choices}
  \choice \texttt{Method}
  \choice \texttt{Field}
  \CorrectChoice \texttt{java.util.List}
  \choice \texttt{Class}
\end{choices}
\begin{solution}
  \texttt{List} lives in \textbf{\texttt{java.util}} and is a collections interface. Reflection code often \emph{returns} ordinary objects or arrays (e.g.\ \texttt{Method[]}), but the \textbf{descriptor types} you navigate with are from \texttt{java.lang.reflect}, not \texttt{java.util}:

\begin{lstlisting}[style=javaexam]
import java.lang.reflect.Method;

Method[] methods = String.class.getDeclaredMethods();
\end{lstlisting}
\end{solution}

\question[4]
What is the Reflection API and what is it used for?
\makeemptybox{1.35in}
\begin{solution}
  Inspecting or instantiating classes, methods, fields, and annotations at runtime; frameworks use it for DI and tooling, with performance and encapsulation trade-offs.
\end{solution}

\newpage

\uplevel{\par\textbf{Threads and concurrency}\par}

\question[4]
What is a thread?
\makeemptybox{1in}
\begin{solution}
  A lightweight unit of execution within a process sharing memory with peers; the JVM maps Java threads to OS threads (platform-dependent).
\end{solution}

\question[4]
How can we create a thread?
\makeemptybox{1.2in}
\begin{solution}
  Subclass \texttt{Thread} and override \texttt{run}, or implement \texttt{Runnable}/\texttt{Callable} and pass to a \texttt{Thread} or executor; prefer executors for production pools.
\end{solution}

\question[4]
How can we implement multithreading in Java?
\makeemptybox{1.25in}
\begin{solution}
  Use threads, \texttt{ExecutorService}, parallel streams, \texttt{CompletableFuture}, or reactive stacks; protect shared mutable state with correct coordination.
\end{solution}

\question[4]
What are some states a thread can be in?
\makeemptybox{1.2in}
\begin{solution}
  E.g.\ \texttt{NEW}, \texttt{RUNNABLE}, blocked or waiting states while parked on locks or \texttt{wait}/\texttt{sleep}, and \texttt{TERMINATED} when \texttt{run} finishes.
\end{solution}

\question[4]
What is synchronization?
\makeemptybox{1.2in}
\begin{solution}
  Coordinating access to shared resources: \texttt{synchronized} blocks/methods, locks, atomics, or concurrent collections to prevent races and ensure visibility.
\end{solution}

\question[4]
What is the difference between deadlock and livelock?
\makeemptybox{1.25in}
\begin{solution}
  Deadlock: threads block forever waiting on each other's locks. Livelock: threads keep changing state in response but make no progress (often ``polite'' retry loops).
\end{solution}

\question[4]
Describe the Producer Consumer problem.
\makeemptybox{1.35in}
\begin{solution}
  Producers enqueue work/items; consumers dequeue; a bounded buffer needs coordination so producers wait when full and consumers when empty---classic use of queues and condition variables/\texttt{BlockingQueue}.
\end{solution}

\newpage

\uplevel{\par\textbf{Threads and concurrency} (\texttt{java.lang.Thread}, \texttt{java.lang.Runnable})\par}

\question[4]
Thread definition. What is a thread?
\begin{choices}
  \choice A special method that executes only once
  \choice A class which allows for more efficient and faster processing
  \CorrectChoice A representation of a path of execution, allowing for concurrent processing
  \choice A separate program which can be run only after the end of the current program
\end{choices}
\begin{solution}
  A \textbf{thread} is a path of execution inside one \textbf{OS process}: its own call stack and program counter, sharing the heap with sibling threads. The class \texttt{java.lang.Thread} is the main API wrapper:

\begin{lstlisting}[style=javaexam]
Thread worker = new Thread(() -> {
    System.out.println("Running on " + Thread.currentThread().getName());
});
worker.start();
\end{lstlisting}
\end{solution}

\question[4]
Threads can be created by extending the \texttt{Runnable} class.
\begin{choices}
  \CorrectChoice False
  \choice True
\end{choices}
\begin{solution}
  \texttt{Runnable} is a \textbf{functional interface} in \texttt{java.lang}---you \textbf{implement} it (or pass a lambda). You \textbf{extend} \texttt{Thread} if you subclass \texttt{Thread} and override \texttt{run()}. Typical pattern:

\begin{lstlisting}[style=javaexam]
Thread t1 = new Thread(() -> { /* task */ });

Thread t2 = new Thread(new Runnable() {
    @Override public void run() { /* task */ }
});
\end{lstlisting}
\end{solution}

\question[4]
The \texttt{run()} method ``comes from'' the \texttt{Thread} class only---it is \textbf{not} specified by the \texttt{Runnable} interface.
\begin{choices}
  \CorrectChoice False
  \choice True
\end{choices}
\begin{solution}
  \texttt{Runnable} \textbf{declares} \texttt{void run()}. \texttt{Thread} \textbf{implements} \texttt{Runnable}. When you use \texttt{new Thread(runnable)}, \texttt{Thread.run()} delegates to that \texttt{Runnable}'s \texttt{run()}:

\begin{lstlisting}[style=javaexam]
Runnable job = () -> System.out.println("job");
new Thread(job).start();
\end{lstlisting}

  So the task contract is tied to \texttt{Runnable}; \texttt{Thread} is the scheduler wrapper and also allows overriding \texttt{run()} on subclasses.
\end{solution}

\question[4]
Which of the following is \textbf{NOT} a \texttt{Thread.State} value in the JVM?
\begin{choices}
  \choice \texttt{RUNNABLE}
  \choice \texttt{BLOCKED}
  \CorrectChoice \texttt{START}
  \choice \texttt{TERMINATED}
  \choice \texttt{WAITING}
\end{choices}
\begin{solution}
  \texttt{Thread.State} is an \texttt{enum} in \texttt{java.lang.Thread}: \texttt{NEW}, \texttt{RUNNABLE}, \texttt{BLOCKED}, \texttt{WAITING}, \texttt{TIMED\_WAITING}, \texttt{TERMINATED}. \texttt{RUNNABLE}, \texttt{BLOCKED}, \texttt{WAITING}, and \texttt{TERMINATED} are all real states. You call \texttt{start()} on a \texttt{NEW} thread; there is no \texttt{START} state.

\begin{lstlisting}[style=javaexam]
Thread t = new Thread(() -> { });
System.out.println(t.getState()); // NEW
t.start();
// Shortly after: usually RUNNABLE, then eventually TERMINATED
\end{lstlisting}
\end{solution}

\question[4]
A thread that is in the \texttt{BLOCKED} state will \emph{immediately} begin running as soon as it becomes unblocked.
\begin{choices}
  \choice TRUE
  \CorrectChoice FALSE
\end{choices}
\begin{solution}
  When a thread leaves \texttt{BLOCKED} (e.g.\ acquired a monitor), it becomes \texttt{RUNNABLE}. The OS / JVM scheduler decides \textbf{when} it actually runs; other runnable threads may run first.

\begin{lstlisting}[style=javaexam]
synchronized (lock) {
    // Other threads blocked on 'lock' become RUNNABLE when released,
    // not guaranteed to run the very next instruction.
}
\end{lstlisting}
\end{solution}

\question[4]
Can we completely avoid deadlock in Java?
\begin{choices}
  \choice TRUE
  \CorrectChoice FALSE
\end{choices}
\begin{solution}
  No general algorithm can prove arbitrary user code deadlock-free. A classic pattern: thread~1 locks \texttt{A} then waits for \texttt{B}; thread~2 locks \texttt{B} then waits for \texttt{A}. Mitigations: fixed global lock order, smaller critical sections, \texttt{tryLock} with timeout, structured concurrency libraries.

\begin{lstlisting}[style=javaexam]
// Illustration only -- bad ordering risk if another thread locks b then a
synchronized (a) {
    synchronized (b) { /* work */ }
}
\end{lstlisting}
\end{solution}

\question[4]
Which of the following method suspend the thread for mentioned time duration in argument?
\begin{choices}
  \choice \texttt{Thread.kill()}
  \choice \texttt{Thread.yield()}
  \choice \texttt{Thread.Suspend()}
  \CorrectChoice \texttt{Thread.sleep()}
\end{choices}
\begin{solution}
  \texttt{Thread.sleep(long millis)} (and overloads with \texttt{TimeUnit} / nanos in newer APIs) pauses the \textbf{current} thread for roughly the requested time without busy-waiting. \texttt{yield()} only hints to the scheduler. \texttt{suspend()} is deprecated. There is no standard \texttt{Thread.kill()}.

\begin{lstlisting}[style=javaexam]
try {
    Thread.sleep(500); // milliseconds
} catch (InterruptedException e) {
    Thread.currentThread().interrupt();
}
\end{lstlisting}
\end{solution}

\endgradingrange{csajavaplus}
